% -*- coding: utf-8 -*-
% !TEX program = xelatex
\documentclass[plain,noanswer,medmath]{../../template/jnuexam}
\usepackage{../../template/kysx}

\begin{document}


\examtitle{2005}{4}


\loadproblems{1}{2005P1}%载入试卷一
\loadproblems{2}{2005P2}%载入试卷二
\loadproblems{3}{2005P3}%载入试卷三


\makepart{填空题}{1～6小题，每小题4分，共24分}


\useproblem{3}{1}{1}%【同试卷三第一(1)题】


\useproblem{3}{1}{2}%【同试卷三第一(2)题】


\useproblem{3}{1}{3}%【同试卷三第一(3)题】


\useproblem{3}{1}{4}%【同试卷三第一(4)题】


\useproblem{1}{1}{5}%【同试卷一第一(5)题】


\useproblem{1}{1}{6}%【同试卷一第一(6)题】


\makepart{选择题}{7～14小题，每小题4分，共32分}


\useproblem{3}{2}{7}%【同试卷三第二(7)题】


\useproblem{3}{2}{8}%【同试卷三第二(8)题】


\begin{problem}
下列结论中正确的是\pickout{}
\begin{abcd}
\item $\int_1^{+\infty} \frac{\dx}{x(x + 1)}$与$\int_0^1 \frac{\dx}{x(x + 1)}$都收敛．
\item $\int_1^{+\infty} \frac{\dx}{x(x + 1)}$与$\int_0^1 \frac{\dx}{x(x + 1)}$都发散．
\item $\int_1^{+\infty} \frac{\dx}{x(x + 1)}$发散，$\int_0^1 \frac{\dx}{x(x + 1)}$收敛．
\item $\int_1^{+\infty} \frac{\dx}{x(x + 1)}$收敛，$\int_0^1 \frac{\dx}{x(x + 1)}$发散．
\end{abcd}
\end{problem}


\useproblem{3}{2}{10}%【同试卷三第二(10)题】


\useproblem{3}{2}{11}%【同试卷三第二(11)题】


\begin{problem}
设$A$, $B$, $C$均为$n$阶矩阵，$E$为$n$阶单位矩阵，若$B = E + AB$, $C = A + CA$，则$B - C$为\pickout{}
\begin{abcd}
\item $E$．
\item $- E$．
\item $A$．
\item $- A$．
\end{abcd}
\end{problem}


\useproblem{1}{2}{13}%【同试卷一第二(13)题】


\begin{problem}
设$X_1,X_2, \cdots ,X_n$为独立同分布的随机变量列，且均服从参数为$\lambda (\lambda > 1)$的指数分布，
记$\Phi (x)$为标准正态分布函数，则\pickout{}
\begin{abcd}
\item $\lim_{n \to \infty} P\Bigg\{\frac{\sum_{i=1}^n X_i - n\lambda}{\lambda\sqrt{n}} \le x \Bigg\}
       = \Phi(x)$．
\item $\lim_{n \to \infty} P\Bigg\{\frac{\sum_{i=1}^n X_i - n\lambda}{\sqrt{n\lambda}} \le x \Bigg\}
       = \Phi(x)$．
\item $\lim_{n \to \infty} P\Bigg\{\frac{\lambda \sum_{i=1}^n X_i - n}{\sqrt{n}} \le x \Bigg\}
       = \Phi(x)$．
\item $\lim_{n \to \infty} P\Bigg\{\frac{\sum_{i=1}^n X_i - \lambda}{\sqrt{n\lambda}} \le x \Bigg\}
       = \Phi(x)$．
\end{abcd}
\end{problem}


\makepart{解答题}{15～23小题，共94分}


\useproblem[points=8]{3}{3}{15}%【同试卷三第三(15)题】


\useproblem[points=8]{3}{3}{16}%【同试卷三第三(16)题】


\useproblem[points=9]{2}{3}{21}%【同试卷二第三(21)题】


%【类似试卷二第三(20)题】
\begin{problem}[points=9]
求$f(x,y) = x^2 - y^2 + 2$在椭圆域$D = \left\{(x,y)\middle|x^2 + \frac{y^2}{4} \le 1 \right\}$
上的最大值和最小值．
\end{problem}


\useproblem[points=8]{3}{3}{19}%【同试卷三第三(19)题】


\useproblem[points=13]{3}{3}{20}%【同试卷三第三(20)题】


\begin{problem}[points=13]
设$A$为三阶矩阵，$\alpha_1,\alpha_2,\alpha_3$是线性无关的三维列向量，且满足
$$A\alpha_1 = \alpha_1 + \alpha_2 + \alpha_3,\quad A\alpha_2 = 2\alpha_2 + \alpha_3,
   \quad A\alpha_3 = 2\alpha_2 + 3\alpha_3.$$
\begin{enumerate}
\item 求矩阵$B$，使得$A(\alpha_1,\alpha_2,\alpha_3) = (\alpha_1,\alpha_2,\alpha_3)B$；
\item 求矩阵$A$的特征值；
\item 求可逆矩阵$P$，使得$P^{-1} AP$为对角矩阵．
\end{enumerate}
\end{problem}


\useproblem[points=13]{3}{3}{22}%【同试卷三第三(22)题】【类似试卷一第三(22)题】


%【类似试卷一第三(23)题】【类似试卷三第三(23)题】
\begin{problem}[points=13]
设$X_1,X_2, \cdots ,X_n(n > 2)$为来自总体$N(0,\sigma^2)$的简单随机样本，
其样本均值为$\overline{X}$，记$Y_i = X_i - \overline{X}$，$i = 1,2, \cdots ,n$．求：
\begin{enumerate}
\item $Y_i$的方差$DY_i$，$i = 1,2, \cdots ,n$；
\item $Y_1$与$Y_n$的协方差$\Cov (Y_1,Y_n)$；
\item $P\{Y_1 + Y_n \le 0\}$．
\end{enumerate}
\end{problem}


\end{document}
